\begin{frame}{Conclusion}
\begin{block}{Main Finding}
Star-LoRA makes AdaLoRA more data-aware and improves over the manually configured AdaLoRA baseline, but fixed LoRA is still the best method on this fake-review dataset.
\end{block}

\vspace{6pt}
\begin{itemize}
  \item Full-data F1: Star-LoRA \(0.9975\), AdaLoRA \(0.9960\), LoRA \(0.9992\).
  \item In the data-size sweep, Star-LoRA is often better than AdaLoRA at very small and very large sample budgets.
  \item The rank policy responds to data size by increasing target rank from 5 to 11 in the sweep, and to 12 for the full run.
\end{itemize}
\end{frame}

\begin{frame}{Limitations}
\begin{itemize}
  \item The current rank policy is hand-designed, so it can over-allocate capacity when a simple fixed-rank LoRA is already sufficient.
  \item Results are from one dataset split and one major task family: fake review classification.
  \item Stability scores are logged and visualized, but they are not yet used directly to update PEFT AdaLoRA's internal allocator.
  \item Star-LoRA has higher runtime than LoRA in the final run because it starts with larger adaptive-rank capacity.
\end{itemize}
\end{frame}

\begin{frame}{Future Work}
\begin{itemize}
  \item Learn the rank policy from multiple datasets instead of using fixed scaling constants.
  \item Add an ablation study for each statistic: size, sparsity, skew, entropy, and embedding variance.
  \item Use stability scores during training, not only as post-hoc logging.
  \item Test whether Star-LoRA helps more on noisier, imbalanced, or lower-resource review datasets.
\end{itemize}
\end{frame}
